Como se ha justificado en el análisis comparativo del estado del arte (\autoref{SS:EA_IA_COMPARATIVA}), optar por un servicio de inferencia local elimina tanto la dependencia económica de APIs comerciales como el envío de datos académicos a terceros. Ollama\footnote{\url{https://ollama.com/}} es la herramienta seleccionada para cumplir este papel, ya que actúa como un servidor local que gestiona la descarga, carga en memoria y ejecución de modelos de lenguaje de código abierto directamente en el hardware del servidor, exponiendo una API REST interna con la que el \textit{backend} se comunica a través de la librería oficial \texttt{ollama-python}\footnote{\url{https://github.com/ollama/ollama-python}}. El sistema está preparado para alternar dinámicamente de modelo sin necesidad de modificar código ni reiniciar el servidor, puesto que la clase \texttt{SystemConfig} almacena el nombre del modelo activo y el panel de administración de Django ofrece un selector dinámico que consulta en tiempo real los modelos disponibles en la instancia de Ollama.